\section{Methodology}
\label{sec:methods}

% ── Overview ──────────────────────────────────────────────────────────────────
\subsection{Overview}

We ask whether the partisan geometry encoded in LLM attention activations tracks real issue partisanship as measured by survey responses. For each of $T$ political topics, we compute two quantities:
\begin{enumerate}
    \item \textbf{Activation-space partisan divergence}: the Mahalanobis distance between Democrat and Republican centroids in a per-head PCA-reduced activation space.
    \item \textbf{Survey partisanship}: the normalized partisan gap in GSS response means.
\end{enumerate}
The main outcome is the Pearson $r$ (and Spearman $\rho$) between these two vectors across all $T$ topics.

% ── GSS ───────────────────────────────────────────────────────────────────────
\subsection{Survey Partisanship Measure}
\label{sec:gss}

For each topic $t$, we compute the survey-measured partisan gap from the General Social Survey (GSS) 2021--2024 cumulative dataset:
\begin{equation}
    \delta_t = \frac{|\bar{x}^R_t - \bar{x}^D_t|}{s_t}
    \label{eq:gss_pol}
\end{equation}
where $\bar{x}^D_t$ and $\bar{x}^R_t$ are the mean response codes among Democrats and Republicans respectively, and $s_t$ is the range of the response scale (i.e., the maximum minus the minimum response code), normalizing $\delta_t \in [0,1]$ regardless of the original scale. Topics are included only if at least 100 Democrats and 100 Republicans responded ($n_D \geq 100$, $n_R \geq 100$, $n \geq 200$ total).

\paragraph{Party coding.}
Respondents are classified by the GSS \texttt{partyid} variable on a 7-point scale:
\begin{center}
\begin{tabular}{clc}
\toprule
\texttt{partyid} & Label & Party \\
\midrule
0 & Strong Democrat           & $D$ \\
1 & Not strong Democrat       & $D$ \\
2 & Independent, near Dem.    & $D$ \\
3 & Independent               & \textit{excluded} \\
4 & Independent, near Rep.    & $R$ \\
5 & Not strong Republican     & $R$ \\
6 & Strong Republican         & $R$ \\
\bottomrule
\end{tabular}
\end{center}

\paragraph{Topic scope.}
We analyze two categories of GSS topics:
\begin{itemize}
    \item \textbf{Public issues} ($T_\text{pub} = 126$): questions about policy, government, and social issues. Eight topics are excluded whose wording contains explicit racial comparison framing that would conflate attitude with world knowledge (\texttt{hubbywk1}, \texttt{racdif1}--\texttt{racdif4}, \texttt{workwhts}, \texttt{wlthwhts}, \texttt{intlwhts}).
    \item \textbf{Private life} ($T_\text{priv} = 73$): questions about personal behavior, religious practice, and lifestyle (e.g., church attendance, sexual attitudes, family structure). Five topics are excluded due to ambiguous framing (\texttt{reborn}, \texttt{marwht}, \texttt{helpful}, \texttt{helpfulnv}, \texttt{helpfulv}).
\end{itemize}

% ── Activation extraction ─────────────────────────────────────────────────────
\subsection{Activation Extraction}
\label{sec:extraction}

We register forward hooks on the output-projection layer (\texttt{o\_proj}) of each of the $L$ transformer layers, capturing \texttt{input[0]} — the pre-projection concatenated head outputs — of shape $[B, S, H{\cdot}D_h]$. This is reshaped to $[B, S, H, D_h]$ and the \textbf{last token position} $[:,{-1},:,:]$ is retained, yielding shape $[B, H, D_h]$.

Capturing the \emph{input} to \texttt{o\_proj} (rather than the output) allows clean decomposition into individual heads, since the projection mixes heads after this point. After all $N$ prompts are processed in batches, activations are concatenated to $\mathbf{X} \in \mathbb{R}^{N \times L \times H \times D_h}$.

% ── PCA Mahalanobis ───────────────────────────────────────────────────────────
\subsection{Per-Head Mahalanobis Distance}
\label{sec:mahal}

For each attention head $(l, h)$, extract $\mathbf{X}_{lh} = \mathbf{X}[:, l, h, :] \in \mathbb{R}^{N \times D_h}$ (typically $D_h = 128$) and proceed as follows:

\begin{enumerate}[label=\arabic*.]
    \item \textbf{PCA}: Fit on all $N$ rows pooled and reduce to $k = \min(15, D_h, N{-}1)$ components, giving $\tilde{\mathbf{X}}_{lh} \in \mathbb{R}^{N \times k}$.

    \item \textbf{Party centroids}: Partition rows into Democrats $\mathcal{D}$ and Republicans $\mathcal{R}$. The centroid for party $P \in \{D, R\}$ is the mean of its members' PCA-projected activations:
    \begin{equation}
        \boldsymbol{\mu}_P = \frac{1}{|\mathcal{P}|}\sum_{i \in \mathcal{P}} \tilde{\mathbf{x}}_i
        \label{eq:centroid}
    \end{equation}
    where $\tilde{\mathbf{x}}_i \in \mathbb{R}^k$ is the PCA-projected activation for subject $i$. As a robustness check, we can also do the coordinate-wise median for the centroid in place of $\boldsymbol{\mu}_P$.

    \item \textbf{Pooled covariance}:
    \begin{equation}
        \boldsymbol{\Sigma} = \tfrac{1}{2}(\boldsymbol{\Sigma}_D + \boldsymbol{\Sigma}_R) + \varepsilon\mathbf{I}, \quad \varepsilon = 10^{-6}
    \end{equation}

    \item \textbf{Mahalanobis distance}:
    \begin{equation}
        m_{lh} = \sqrt{(\boldsymbol{\mu}_D - \boldsymbol{\mu}_R)^\top \boldsymbol{\Sigma}^{-1} (\boldsymbol{\mu}_D - \boldsymbol{\mu}_R)}
        \label{eq:mahal}
    \end{equation}
\end{enumerate}

This is computed in parallel for each of the $L \times H$ heads, producing a grid $\mathbf{M} \in \mathbb{R}^{L \times H}$. Two summary statistics aggregate the grid:
\begin{align}
    m^\text{all} &= \frac{1}{LH}\sum_{l,h} m_{lh} \label{eq:mall} \\
    m^\text{mid} &= \text{mean}_{l \in \mathcal{L}_\text{mid},\, h}\; m_{lh} \label{eq:mmid}
\end{align}
where $\mathcal{L}_\text{mid} = \{l : \lfloor 0.45L \rfloor \le l < \max(\lfloor 0.55L \rfloor, \lfloor 0.45L \rfloor{+}1)\}$ selects the middle 10\% of layers. Empirically, $m^\text{mid}$ performs worse than $m^\text{all}$ for deep models ($L \ge 70$), with correlations reduced by up to 0.48.

% ── Simulation methods ────────────────────────────────────────────────────────
\subsection{Simulation Methods}
\label{sec:sim}

We test two methods for constructing the $N$ prompts per topic.

\subsubsection{Politician Simulation}
\label{sec:politician}

\paragraph{Subjects.} 550 members of the 116th U.S. Congress with valid name and party affiliation from the DW-NOMINATE dataset.

\paragraph{Prompt conditions.} One prompt is generated per politician per topic ($N = 550$). Each prompt is prefixed by a system message that frames the politician-simulation task. The system message was varied across nine phrasings to verify robustness; alignment differs by less than $0.025$ across all variants (Appendix~\ref{app:sys_msg}). Three user-turn conditions are tested:

\begin{center}
\small
\begin{tabular}{lp{5.5cm}p{5.5cm}}
\toprule
Condition & Instruct template & Base template \\
\midrule
\textsc{rhetorical} & \texttt{Generate a statement by \{name\} on \{topic\}.} & \texttt{\{name\} makes a statement on \{topic\}:} \\
\textsc{stance}     & \texttt{What is \{name\}'s position on \{topic\}?}     & \texttt{On \{topic\}, \{name\}'s position is} \\
\textsc{survey}     & \texttt{If asked in a national survey about \{topic\}, how would \{name\} respond?} & \texttt{In a survey about \{topic\}, \{name\} would respond} \\
\bottomrule
\end{tabular}
\end{center}

\texttt{\{topic\}} is a natural-language clause derived from the GSS question wording. Table~\ref{tab:topic_examples} lists representative topics from both categories. Public-issue topics cover policy and governance; private-life topics cover personal behaviour, moral attitudes, and lifestyle---areas where partisan sorting is less explicit in political discourse but nonetheless present in survey data.

\begin{table}[ht]
\centering
\small
\caption{Sample GSS topics and their natural-language clauses used in prompts.}
\label{tab:topic_examples}
\begin{tabular}{llp{8cm}}
\toprule
Variable & Category & Natural-language clause \\
\midrule
\texttt{abhelp2}  & public  & whether to help pay for an abortion \\
\texttt{gunlaw}   & public  & whether to favor or oppose requiring a permit to buy a gun \\
\texttt{eqwlth}   & public  & whether the government should reduce income differences between the rich and poor \\
\texttt{grass}    & public  & whether marijuana should be made legal \\
\texttt{confinan} & public  & the level of confidence in banks and financial institutions \\
\texttt{nextgen}  & public  & whether science and technology will give more opportunities to the next generation \\
\midrule
\texttt{spanking}  & private & whether it is sometimes necessary to discipline a child with spanking \\
\texttt{evstray}   & private & whether to have sex with someone other than a spouse while married \\
\texttt{marmakid}  & private & whether a single mother can raise a child as well as a married couple \\
\texttt{fear}      & private & whether to be afraid to walk alone at night in the neighborhood \\
\texttt{life}      & private & whether life is exciting, routine, or dull \\
\texttt{marasian}  & private & how to feel about a close relative marrying an Asian American person \\
\bottomrule
\end{tabular}
\end{table}

To illustrate the full prompts fed to the model, the following are concrete instantiations across prompt conditions and topics, each preceded by the system message \textit{``You are simulating the public stance of U.S.\ politicians.''}:

\begin{itemize}[leftmargin=*]
    \item \textsc{stance}, public topic: \textit{``What is Ted Cruz's position on whether to help pay for an abortion?''}
    \item \textsc{rhetorical}, public topic: \textit{``Generate a statement by Alexandria Ocasio-Cortez on the level of confidence in banks and financial institutions.''}
    \item \textsc{survey}, public topic: \textit{``If asked in a national survey about whether to favor or oppose requiring a permit to buy a gun, how would Mitt Romney respond?''}
    \item \textsc{stance}, public topic: \textit{``What is Elizabeth Warren's position on whether marijuana should be made legal?''}
    \item \textsc{survey}, private topic: \textit{``If asked in a national survey about whether to have sex with someone other than a spouse while married, how would Mitt Romney respond?''}
    \item \textsc{stance}, private topic: \textit{``What is Nancy Pelosi's position on whether it is sometimes necessary to discipline a child with spanking?''}
    \item \textsc{rhetorical}, private topic: \textit{``Generate a statement by Marsha Blackburn on whether a single mother can raise a child as well as a married couple.''}
    \item \textsc{survey}, private topic: \textit{``If asked in a national survey about whether life is exciting, routine, or dull, how would Marco Rubio respond?''}
\end{itemize}

The private-life examples highlight a distinctive feature of this study: models are asked to simulate politicians on topics---marital fidelity, parenting practices, neighbourhood safety---that rarely appear in legislative records or campaign rhetoric, yet still exhibit partisan structure in survey data.

The three conditions test distinct representational hypotheses: \textsc{rhetorical} tests whether political speech generation is organized by partisan identity; \textsc{stance} tests factual belief-attribution; \textsc{survey} tests whether the model tracks response-style differences between parties.

\subsubsection{Demographic Simulation}
\label{sec:demo}

\paragraph{Subjects.} For each topic $t$, we include all GSS 2021–2024 respondents who (a) provided a valid response to topic $t$ and (b) are classified as Democrat or Republican. We evaluate two sampling settings:
(1) a \textbf{stratified sampling} condition, in which a 10\% sample is drawn stratified on \texttt{polviews}, \texttt{age\_bin}, \texttt{degree}, \texttt{race}, \texttt{sex}, and \texttt{rincome}, requiring at least 10 respondents per party; and
(2) a \textbf{non-stratified} condition, where respondents are sampled without demographic stratification. This allows us to assess the impact of sampling design on simulation outcomes.

\paragraph{Persona profile.} Each respondent is represented by up to 83 demographic fields. For each non-missing field, numeric codes are mapped to human-readable strings using Stata value labels (e.g., \textit{`Age: 45''}, \textit{`Education: Bachelor’s degree''}). Fields are concatenated using ``\texttt{. }'' as a separator and \textbf{randomly shuffled per topic} (seeded by \texttt{hash(topic)}) to mitigate ordering artifacts.

In addition to the full-profile setting, we conduct a \textbf{focused demographic study} using a reduced profile containing only key attributes: \texttt{age}, \texttt{sex}, \texttt{race}, \texttt{relig}, and \texttt{occ}. This enables analysis of how model behavior changes when only core demographic signals are provided.

\paragraph{Prompt (Format B, default).}
\begin{quote}
\textit{System}: ``You are simulating the views of an American.''\\[2pt]
\textit{User}: ``Given the following background about a person:\\
\quad\texttt{\{profile\}}\\[2pt]
How would they answer: \texttt{\{survey\_question\}}''
\end{quote}

To prevent leakage, any lifestyle variable that coincides with the current survey topic is excluded from that respondent's profile for that topic.

% ── Alignment metric ──────────────────────────────────────────────────────────
\subsection{Alignment Metric}
\label{sec:alignment}

For a given simulation method, condition, and layer-selection variant, we correlate the vector of activation-space partisan divergence scores $\hat{\boldsymbol{\delta}}$ with the GSS partisanship vector $\boldsymbol{\delta}$ across all $T$ topics in the category ($T_\text{pub} = 126$ or $T_\text{priv} = 73$):
\begin{align}
    r &= \mathrm{Pearson}(\hat{\boldsymbol{\delta}},\, \boldsymbol{\delta}) \\
    \rho &= \mathrm{Spearman}(\hat{\boldsymbol{\delta}},\, \boldsymbol{\delta})
\end{align}

A high $r$ indicates that topics generating high partisan divergence in LLM activations are the same topics where actual survey respondents diverge most.

\paragraph{Politician vs.\ demographic comparison.}
When $r_\text{politician} > r_\text{demo}$, the model's parametric political knowledge (accessed via named politician prompts) more faithfully mirrors the GSS partisanship landscape than explicit demographic conditioning on real respondent data. This reflects a difference between what the model \emph{knows} about which topics are contested between parties and how it \emph{behaves} when role-playing constrained personas.

% ── Cultural Terrain Discovery ────────────────────────────────────────────────
\subsection{Cultural Terrain Discovery}
\label{sec:discovery}

The preceding analyses validate the activation-based partisan-divergence measure against known GSS topics. We additionally apply the same measure in an \emph{unsupervised discovery} mode: given a large corpus of novel opinion clauses not drawn from any existing survey instrument, which topics does the model flag as highly partisan in activation space? High-scoring topics from this scan are candidates for survey items that existing instruments do not cover.

\paragraph{Corpus construction.}
We assembled a corpus of approximately 1.3 million candidate clauses from four source collections (Table~\ref{tab:discovery_corpora}). Raw items from each source --- bill titles, academic abstracts, social-media posts, and encyclopaedic descriptions --- were converted into natural-language opinion clauses of the form ``the [noun phrase]'' or ``whether [subordinate clause]'', matching the format of GSS topic wording. Clause generation used an instruction-tuned model prompted to reframe each item as a topic on which people can hold opposing views; items that could not be meaningfully reformulated (purely factual content, data artefacts) were discarded.

\begin{table}[ht]
\centering
\small
\caption{Source corpora for cultural terrain discovery.}
\label{tab:discovery_corpora}
\begin{tabular}{lrp{7.5cm}}
\toprule
Source & Clauses & Content \\
\midrule
Merged (multi-source) & 477,343 & Wikipedia articles, Reddit opinion posts, news headlines, memes, party platforms, ballot measures, book and film titles, ATUS time-use labels \\
Congress bills        & 359,302 & U.S. Congress bill titles and CRS report summaries \\
Web of Science        & 424,445 & Academic abstract text, restricted to social-science and humanities journals \\
Reddit (partisan)     &  47,018 & Posts from politically-aligned subreddits, filtered to surface-apolitical wording \\
\midrule
Total                 & 1,308,108 & \\
\bottomrule
\end{tabular}
\end{table}

\paragraph{Scoring.}
Each clause is scored using the identical pipeline as Section~\ref{sec:mahal}: the same 550 Congress members, the same \textsc{stance} prompt template, the same \texttt{o\_proj} activation extraction, PCA reduction ($k = 15$), and mean Mahalanobis distance $m^\text{all}$ averaged across all $L \times H$ heads. The scoring model is dolphin-2.9.3-llama-3-8b (8B) with the roleplay system message (Appendix~\ref{app:sys_msg}). Because each clause is scored independently using the same partisan geometry that correlates $r \approx 0.74$ with GSS partisanship on known topics, the $m^\text{all}$ score serves as a proxy for expected survey-measured partisanship.

\paragraph{Output.}
The top-scoring clauses from each corpus are retained as candidate survey items. Table~\ref{tab:discovery_top} (Section~\ref{sec:results_discovery}) lists representative high-scoring examples. Because the scoring is unsupervised --- no ground-truth partisanship labels are used --- the discovered topics represent the model's implicit map of the partisan cultural terrain, extending well beyond the policy and lifestyle domains covered by the GSS.

\paragraph{Surprise analysis: latent partisan structure.}
A high $m^\text{all}$ score is not itself surprising if the clause is overtly political --- a bill about abortion policy or an article titled ``Republican efforts to restrict voting'' would be expected to produce partisan activation geometry. The more scientifically interesting case is a topic that is \emph{more} polarized than its surface content predicts. We isolate these cases via a two-stage regression procedure.

\textit{Stage 1 --- Political content scoring.} Each clause is embedded with a 300M-parameter Gemma3-backbone sentence encoder (mean-pooled, L2-normalised, 768-dimensional). The political content of a clause is operationalised as its cosine similarity to a set of 234 Ballotpedia ballot-measure subtopics spanning ten parent categories (Budget and Tax, Criminal Justice, Education, Elections, Environment, Governance, Infrastructure, Labour and Healthcare, Military, and Social Issues). These anchor similarities form a feature vector of political salience.

\textit{Stage 2 --- Per-corpus OLS and residual.} Within each corpus separately, we regress $m^\text{all}$ on the 234 anchor-similarity features using ordinary least squares with 5-fold cross-validation ($K = 5$). Fitting per-corpus rather than pooled avoids confounding the baseline polarization level across sources (congressional bill clauses are almost universally political; AITA interpersonal-dilemma posts are almost universally apolitical). The surprise score for clause $i$ in corpus $c$ is the OLS residual:
\begin{equation}
    s_i = m^\text{all}_i - \hat{m}^\text{all}_i
\end{equation}
where $\hat{m}^\text{all}_i$ is the within-corpus predicted value. A large positive residual indicates that the model's activation geometry flags the topic as politically contentious even though the clause's semantic content offers little indication of political relevance.

\textit{Non-political filter.} To surface genuinely surprising topics rather than noisily phrased political ones, we restrict attention to clauses whose predicted polarization $\hat{m}^\text{all}_i$ falls below the corpus mean. Within this filtered set, clauses are ranked by $s_i$ descending. We retain the top 1,000 per eligible corpus (AITA and Reddit; WoS, Wikipedia, and Congress are too uniformly political or too sparse after filtering).

\textit{Out-of-sample validation.} To test whether the surprise score generalises, we hand-curated 21 contemporary topics not drawn from any scraped corpus --- including items such as the tradwife lifestyle, raw milk consumption, fluoride removal from drinking water, and similar culturally-charged but explicitly non-partisan issues. We first predict their polarization using the per-corpus Ridge model, then measure actual $m^\text{all}$ scores using the dolphin-2.9.3-llama-3-8b pipeline, and compare predicted versus actual to confirm that the surprise decomposition transfers to genuinely novel clauses.

% ── Question Fundamentalness ──────────────────────────────────────────────────
\subsection{Question Fundamentalness (will elaborate later)}
\label{sec:fundamentalness}

A complementary analysis asks which survey topics are structurally \emph{central} to the opinion landscape --- topics whose activation geometry is most informative and whose responses are most predictive of others.

\paragraph{LLM-side measure.}
For each topic $t$, the per-head Mahalanobis scores $\{m_{lh}\}$ (Section~\ref{sec:mahal}) are aggregated into a single activation-based fundamentalness score $m^\text{all}_t$ (Eq.~\eqref{eq:mall}), averaged across all $L \times H$ heads. A high score indicates that the model's attention heads reliably separate partisan clusters for that topic.

\paragraph{Survey-side measures.}
We evaluate five independent structural measures computed from GSS response data (11,066 respondents, 98 questions appearing in at least two of the three survey years 2021, 2022, 2024):
\begin{enumerate}[label=(\arabic*)]
    \item \textbf{Mutual information}: normalized mutual information (NMI) between each q (uestion pair; a question's fundamentalness score is the sum of its pairwise NMI across all other questions.
    \item \textbf{Predictive power}: out-of-sample $R^2$ when predicting each question's responses from all others; more predictable questions are scored as more fundamental.
    \item \textbf{Network centrality}: a question-correlation graph is constructed with edge weights equal to absolute Pearson $r$ (could also use MI here,  not sure which one is better); fundamentalness is measured by degree, betweenness, closeness, eigenvector, and related centrality variants across graph thresholds.
    \item \textbf{Tree structure}: a Chow-Liu algorithm fits an optimal dependency tree on pairwise MI; fundamentalness is proximity to the tree root (inverse depth, subtree size).
\end{enumerate}
The main outcome is the Pearson correlation between the LLM activation score $m^\text{all}_t$ and each survey-side fundamentalness measure across the shared topics.
